\documentclass{beamer}

\usepackage[utf8]{inputenc}
\usepackage{url}
\usepackage{listings}
\usepackage{hyperref}
\usepackage{xcolor}
\usepackage{graphicx}

\input{../helper}


\title{Introdução à Computação \\ números hexadecimais}
\author{Alexandre Rademaker}
\date{}

\begin{document}

\begin{frame}
 \maketitle
\end{frame}


\begin{frame}[fragile,allowframebreaks]{números hexadecimais}

  A maior parte das culturas ocidentais usa o sistema decimal
  (base-10) para contagem:

  \[ 0 \quad 1 \quad 2 \quad 3 \quad 4 \quad 5 \quad 6 \quad 7 \quad 8 \quad 9 \quad 10 \]

  Como já vimos, os computadores usam o sistema binário, base-2, para
  representar dados.

  \[ 0 \quad 1 \]

  \framebreak

  Como programadores, seria conveniente representarmos e usarmos os
  dados em binário, como os computadores.

  No entanto, entender uma cadeia muito longa de 0s e 1s pode ser bem
  complicado.

  \url{https://en.wikipedia.org/wiki/Hexadecimal}

  \framebreak

  A representação em hexadecimais é bem mais concisa para representar
  dados.

  \[
    0 \quad 1 \quad 2 \quad 3 \quad 4 \quad 5 \quad 6 \quad 7 \quad 8 \quad 9 \quad A \quad B \quad C \quad D \quad E \quad F 
  \]

  O mapeamento de hexadecimais para binários é simples. Um grupo de 4
  bits é capaz de representar 16 diferentes combinações, cada uma
  delas mapeada para um único hexadecimal.

  \framebreak

  \begin{columns}
    \begin{column}{.5\textwidth}
      \centering
      \begin{tabular}{ccc}
        \hline
        dec & bin & hex \\
        \hline
        0 & 0000 & 0\\
        1 & 0001 & 1\\
        2 & 0010 & 2\\
        3 & 0011 & 3\\
        4 & 0100 & 4\\
        5 & 0101 & 5\\
        6 & 0110 & 6\\
        7 & 0111 & 7\\
        \hline
      \end{tabular}
    \end{column}
    \begin{column}{.5\textwidth}
      \centering
      \begin{tabular}{ccc}
        \hline
        dec & bin & hex \\
        \hline
        8 & 1000 & 8\\
        9 & 1001 & 9\\
        10 & 1010 & A\\
        11 & 1011 & B\\
        12 & 1100 & C\\
        13 & 1101 & D\\
        14 & 1110 & E\\
        15 & 1111 & F\\
        \hline
      \end{tabular}
    \end{column}
  \end{columns}

  \framebreak

  \begin{columns}
    \begin{column}{.5\textwidth}
      \centering
      \begin{tabular}{ccc}
        \hline
        dec & bin & hex\\
        \hline
        0 & 0000 & {\color{orange} 0x}0\\
        1 & 0001 & {\color{orange} 0x}1\\
        2 & 0010 & {\color{orange} 0x}2\\
        3 & 0011 & {\color{orange} 0x}3\\
        4 & 0100 & {\color{orange} 0x}4\\
        5 & 0101 & {\color{orange} 0x}5\\
        6 & 0110 & {\color{orange} 0x}6\\
        7 & 0111 & {\color{orange} 0x}7\\
        \hline
      \end{tabular}
    \end{column}
    \begin{column}{.5\textwidth}
      \centering
      \begin{tabular}{ccc}
        \hline
        dec & bin & hex\\
        \hline
        8 & 1000 & {\color{orange} 0x}8\\
        9 & 1001 & {\color{orange} 0x}9\\
        10 & 1010 & {\color{orange} 0x}A\\
        11 & 1011 & {\color{orange} 0x}B\\
        12 & 1100 & {\color{orange} 0x}C\\
        13 & 1101 & {\color{orange} 0x}D\\
        14 & 1110 & {\color{orange} 0x}E\\
        15 & 1111 & {\color{orange} 0x}F\\
        \hline
      \end{tabular}
    \end{column}
  \end{columns}

  \framebreak

  Assim como ocorre no sistema binário ($2^0$, $2^1$, $2^2$) e decimal
  ($10^0$, $10^1$, $10^2$, $10^3$, etc), a posição do algarismo no
  sistema hexadecimal também altera seu valor.

  \begin{center}
    \renewcommand{\arraystretch}{1.5}
    \begin{tabular}{|l|c|c|c|}
      \hline
      & $16^2$ & $16^1$ & $16^0$ \\
      \hline
      & 256 & 16 & 1\\
      \hline
      0x & 3 & 9 & 7\\
      \hline
    \end{tabular}
  \end{center}

  $397_{16} = (3 \times 16^2) + (9 \times 16^1) + (7 \times 16^0) = 919_{10}$

  \framebreak

  Mas como converter binários em hexadecimais?

  Vamos agrupar os bits em grupos de 4, da direita para a esquerda,
  adicionado zeros à esquerda se necessário.

  Então usamos a tabela de conversão apresentada nos slides
  anteriores.

  \framebreak

  \[ 01000110101000101011100100111101 \]

  \[ 0100 \quad 0110 \quad 1010 \quad 0010 \quad 1011 \quad 1001 \quad 0011 \quad 1101 \]

  \[ 4 \qquad 6 \qquad A \qquad 2 \qquad B \qquad 9 \qquad 3 \qquad D \]

  \[ 0x46A2B93D \]

\end{frame}


\end{document}

%%% Local Variables:
%%% mode: latex
%%% TeX-master: t
%%% End:
